\documentclass[book.tex]{subfiles}

\begin{document}
\chapter{Gauge Equivariant Normalizing Flows for Lattice Field Theory}

\label{chap:gauge_equivariant_flow}
\noindent\rule{11cm}{0.4pt} \\
\textbf{Prerequisites:} \autoref{chap:normalizing_flows}, \autoref{chap:qft_basics} \\
\textbf{Difficulty Level:} ** \\
\noindent\rule{11cm}{0.4pt}
\vspace{5mm}

%"Monte Carlo sampling is a powerful approach to probing quantum field theories defined on a discrete spacetime lattice ("lattice field theories"). We will highlight two recent papers detailing the use of normalizing flows to generate such samples without compromising the exactness of calculations. Gauge symmetry plays a key role in quantum field theories relevant for the Standard Model of particle physics, and our work lays out a framework for exactly encoding such gauge symmetries into these models, making training practically viable. We have shown that normalizing flow models successfully learn to sample distributions relevant to simple lattice field theories with gauge symmetries, and we will demonstrate constructing and training models for these theories in the code walk-through."

Monte Carlo methods allow for sampling quantum field configurations defined on a discrete space time lattice. Such discrete configurations are broadly termed lattice field theories. Unfortunately performing this Monte Carlo sampling can in general be very complex. Normalizing flows can facilitate sampling of these complex distributions. Critically, these normalizing flows must respect gauge symmetries, which are fundamental for the quantum field theories used in the standard model of particle physics. In this chapter, we study a framework to encode gauge symmetries into normalizing flows and discuss how to train the resulting networks \cite{kanwar2020equivariant}, \cite{albergo2019flow}, \cite{boyda2021sampling}. The resulting gauge equivariant normalizing flows can successfully sample distributions on simple lattice field theories with gauge symmetries.

\section{A Review of Lattice Field Theory}

Quantum field theory provides a structure that unifies relativistic quantum mechanics with the electromagnetic, weak, and strong forces. As we have learned earlier, a challenge of the quantum field framework is that we must construct fields of operators spread over all of spacetime, a computationally challenging tasks.

Lattice quantum field theories discretize spacetime as a grid. This discretization allows for computational modeling of quantum field theory with computation systems. Suppose we have an action $S$ defined as below. We can derive a probability density $p$ from this action

\begin{align}
S(\phi) &= \sum_x \sum_y \phi(x) \square(x, y) \phi(y) + m^2 \phi(x)^2 + \phi(y)^4 \\
p(\phi) &= e^{-S(\phi)}/Z
\end{align}


Lattice configurations can be sampled according to this energy function. Deriving these samples has traditionally required sophisticate Markov Chain Monte Carlo (MCMC) samplers. However, there are critical limitations to MCMC style methods. In particular, mixing rates can be very slow, requiring extensive computation. Reweighting techniques can speed up the convergence rate, but it remains a formidable challenge.


Roughly, the sampling process is analogous to that for sampling images. There is a few differences though; the exact probability distribution $p(\phi)$ can be computed. Training datasets are also much more limited, derived from expensive MCMC calculations, but typical lattice sizes can be very large.%You also end up having small training sets of roughly $10^3$ configurations, but each configuration has between $10^9$ and $10^{12}$ numbers. 

Normalizing flows as we have seen previously in the book are a powerful tool for sampling probability distributions. We shall show how to use normalizing flows to sample these complex probability distributions. The normalizing flow acts as a variational ansatz for the function $q$ that we sample from. By normalizing this function properly, we can make the sample close to the true probability distribution.

\begin{align}
 q(\phi) &= r(z) \left | \det_{ij} \frac{\partial [f(z)]_j]}{\partial z_j} \right |^{-1}
\end{align}
This last term is the log-det-Jacobian (LDJ). Here $f$ is a composition of many constituent functions $g_i$
\begin{align}
f = g_n \cdot g_{n-1} \cdots \cdot g_1
\end{align}
If we construct each $g$ to act on only a subset of the lattice components, the log-det-Jacobian becomes upper triangular. For scalar field theory, we can think of this construction as building a checkerboard mask on the lattice. If the checkerboard is black and white, we can think of white squares as frozen and the black squares are updated. 

We also introduce context functions, which are convolutional networks that take the local parameters of the network. These functions take in the frozen sites. The normalizing flow $q$ can be trained by computing the KL divergence between $q$ and the true distribution $p$

\begin{align}
\mathcal{D}_{q||p} &= \mathbb{E}_{\phi \sim q}[\log(q(\phi)) - \log(p(\phi))]
\end{align}

Samples $p$ are drawn from the true distribution and arecomputed by simulation. The normalizing flow $q$ is trained to minimize $\mathcal{D}_{q||p}$

\section{Lattice Gauge Theory}

Quantum chromodynamics is the quantum field theory of the strong nuclear force. Gauge fields are the force carriers for the strong nuclear force. A gauge symmetry is roughly a coordinated local change of variables. 

Matter fields live at the vertices of the lattice and gauge fields live on the edges.

\begin{align}
U_{\mu}(x) \in \textrm{SU}(N)
\end{align} 

Here $\textrm{SU}(N)$ is called the Gauge group.

\begin{align}
    S(U) &= -\frac{\beta}{N} \sum_x \mathrm{Re}\mathrm{tr}[P_{01}(x)] \\
    P_{\mu\nu}(x) &= U_{\mu}(x)U_{\nu}(x + \hat{mu})U_{mu}^{\dagger}(x + \hat{\nu})U_{\nu}^{\dagger}(x)
\end{align}
This last equation is called a Wilson loop.
Lattice gauge theories have a number of symmetries, including discrete translational symmetry, hypercubic symmetry, and guage symmetry. Exploiting these symmetries can yield data efficiency in training.
% (\textcolor{red}{TODO: what is this})

We can impose these symmetries by imposing invariance to the symmetry in the prior and imposing equivariance on the flow $f$. Translational equivariance for example can be achieved by using convolutional networks in context functions and symmetric masking patterns. 

Gauge equivariance is trickier. The intuition is that we want to act on gauge invariant quantities only. Gauge fixing with an explicit factorization doesn't preserve translational symmetry. Gauge fixing with an implicit factorization is hard to preserve across the coupling layer. The solution that is reached is to transform group-valued untraced Wilson loops, and absorb updates using link representations.
\begin{align}
P_{\mu\nu}' &= h(P_{\mu\nu}|I) \\
U_{\mu}' &= P'_{\mu\nu}(P')_{\mu\nu}^{\dagger}U_{\mu} %\\
%   P'_{\mu\nu} 
\end{align}

The core idea is to construct an equivariant kernel while keeping the form of the log-det-Jacobian tractable. We do this by enforcing equivariance under matrix conjugation.
%\begin{enumerate}
%    \item Invertible and tractable LDJ
%    \item Euivariant under matrix conjugaation
\begin{align}
h(XPX^{-1}) = Xh(p)X^{-1}
\end{align}
%    \item \textcolor{red}{TODO}
%\end{enumerate}

Let's take a look at $U(1)$ gauge theory. In this case, the "matrices" are just scalars. The flow based method seems to agree well with the known analytical calculations. This can also be done for $\textrm{SU}(N)$. You want to look at the Haar measure to implement the $\textrm{SU}(N)$-invariant kernel for the flow. So far, this has been done up to $\textrm{SU}(3)$, but techniques still need to be worked out for higher dimensions. 
\end{document}